\section{工作总结}
本文完成了基于已实现 QCS 工程的原理、方案、结果与应用整理。
检测引擎通过源码、符号、字节、已知函数和运算语义通道，
将源码、程序及受支持固件组织为可追溯证据；
国密模块扩充 SM2、SM3、SM4，可选 Miasm 提供有限 CFG、保守合并及带输入身份的摘要。
报告区分接口引用、实现包含与运行时使用，
保留未知配置、共享结构、预算耗尽与失败状态。

本轮将算法主实验扩展至理论章节的全部十三种核心算法，
从开发总体全取源码与 x86-64/ARM64 程序，共 1071 份不同内容输入实际新扫，
TP/FP/FN 为 700/0/47，总体实现召回率 93.71\%、漏报率 6.29\%。
逐算法的漏报目标为不超过 10\%，表格保留绝对数、达标状态和分载体差异。
该结果是已知家族回归，原五算法筛选混合演示、
完整开发、历史覆盖和 Miasm 收益对照继续保留，支持进一步复核。
作品优势仍是证据可审查、多载体接入与资源可控，
未来需要用新的独立实现验证更广范围的覆盖。

GridLens 把识别结果落到虚拟电力边缘网关的发布前体检：
四组件、两版本、两架构和六文件范围形成可演示闭环，
遗留摘要变化有真实源码及扫描证据，公钥迁移保持人工待办。
离线网页与本地服务使演示不依赖实体硬件；这一可用性成果与泛化指标分别评价。
另补齐五类电力终端的软件级场景：不同业务程序与独立清单构成 10 份双架构自制固件，
实际解包核验 80 个文件并扫描 20 个 ELF；22 个正标签全部检出，保留一项误报。
五类 x86-64 业务检查通过，ARM64 只完成静态验证；
这些证据支持固件软件流程，不等同于真实厂商或实体终端覆盖。

\section{后续改进路线}
\subsection{优先提升未知实现覆盖}
首先针对大数与曲线算法建立带 ABI、参数身份、输出缓冲区、成功路径和调用副作用的摘要。
由模乘、Montgomery 约减、模逆、模幂和点运算关系逐步组合成算法见证，
避免仅凭乘法或点乘命名协议。
对 SHA-256、SM3、SM4 继续扩大开发实现多样性，
验证轮常量使用、动态查表、位切片、硬件指令和优化下的数据依赖。
新规则必须先在开发范围完成正反例与预算检验。

\subsection{完善语义与资源边界}
继续处理 Thumb IT、MIPS 延迟槽、调用摘要及别名内存输出，
改进候选选择和启动成本，并检查预算不足时的稳定弃权。
原语或部分表达式不能被包装成整个函数恢复，
也不把生成伪代码后再做名称扫描作为覆盖提升的充分证据。
既有 ARM 探针异常和预算相关回归需要单独修复，再运行相应验证。

\subsection{提高评测的外部有效性}
重新冻结后保留未见作者与构建链的独立留出集合，
按家族汇总并报告各算法、架构、载体和不完整状态。
完整开发、历史观察与演示筛选继续分离，
不在看到留出结果后挑选最有利变体再称其为独立性能。
如需比较其他工具或模型，应固定同一真值、输出范围、预算和计数单位，
不能用各自论文的数字拼接排行榜。

\subsection{从迁移辅助走向工程验证}
在应用层补充配置证据、证书链、依赖版本与标准 CBOM 导出，
再验证真正的网关进程、升级签名、防回滚、消息认证和密钥生命周期。
后量子或混合协议应经过测试向量、互操作、资源开销及接收端兼容实验，
完成这些证据后才能把候选路线改为已验证迁移结果。
Windows 原生与 WSL 的安装演示、真实设备或受控模拟环境也应单独验收。

\section{结语}
发布前密码资产体检的首要价值，是把不可见的密码依赖变成可定位、可解释、
可复查的对象，并为迁移决策保留明确证据。
本作品已经实现这一软件原型及应用展示；未知实现识别与端到端迁移验证仍是下一阶段的核心工作。
